\subsection{Exact contractions of terminal ruled and plane components}

\begin{statement}{C.2.UN.64}{Proved}{Terminal ruled pairs, their actual contractions and the full normal extension}
Let $f:U\to\Delta$ be a proper degeneration of connected complex
two-tori from a smooth threefold. Assume that the full central fibre
$S$ is reduced and SNC. Let $T$ be one component, meeting the rest
of $S$ in exactly one smooth connected curve $C$, contained in a
single other component $P$. Let $k_T,k_P$ be the orders of the
canonical evaluation along these components. Suppose that either
$q:T\to B$ is a specified holomorphic $\PP^1$-bundle over a smooth
compact curve, or $T\simeq\PP^2$.

For a ruled $T$, precisely one of the following two numerical and
marked-pair descriptions holds. Either $C$ is a section,
\[
 K_T=-2C,\qquad k_T-k_P=1,\qquad
 N_{C/T}\simeq K_B^{-1},\qquad N_{C/P}\simeq K_B,
                                                        \tag{C.EL.1}
\]
where the section identifies $C$ with $B$; or $B\simeq\PP^1$,
$C$ is a smooth anticanonical elliptic bisection and
\[
 K_T=-C,\qquad k_T=k_P,\qquad
 T\simeq\mathbb F_e\quad\hbox{with }e\in\{0,1,2\}.
                                                        \tag{C.EL.2}
\]
For $T\simeq\PP^2$, either $C$ is a line and $k_T-k_P=2$,
or $C$ is a smooth cubic and $k_T=k_P$.
These assertions classify the stated terminal pairs in actual
degenerations; they make no occurrence assertion for the
anticanonical alternatives.

In the section case (C.EL.1) there is an actual proper holomorphic
map $\pi:U\to U'$ onto a smooth threefold, equal to an isomorphism
off $T$, which contracts the specified ruling to a smooth curve
$B\subset U'$. It identifies $U$ with the ordinary blowup of $U'$
along $B$. The map $f$ descends to a proper holomorphic torus
degeneration $f':U'\to\Delta$, smooth along $B$; the restriction
$P\to P'$ is an isomorphism and identifies $C$ with $B$.
The full new fibre $S'$ is again reduced and SNC. Its exact normal
bundle and Euler comparisons are
\[
\begin{gathered}
 0\longrightarrow K_B\longrightarrow V:=N_{B/U'}
        \xrightarrow{\,df'\,}\OO_B\longrightarrow0,\qquad
 (T,C)\simeq\bigl(\PP_B(V),\PP_B(K_B)\bigr),\\
 N_{T/U}\simeq\OO_{\PP_B(V)}(-1)\simeq\OO_T(-C),\qquad
 e(S)=e(S')+2-2g(B).
\end{gathered}                                                   \tag{C.EL.3}
\]
Here $\PP_B(V)$ parametrizes lines. The extension class belongs
to the actual one-dimensional space $H^1(B,K_B)$ and is retained,
with its endpoint identifications, by the contraction record.

In the plane-line case the corresponding map contracts $T$ to a
smooth point of a smooth reduced central component, identifies
$U$ with its point blowup, and gives
\[
       P=\operatorname{Bl}_{p}P',\qquad e(S)=e(S')+2.
                                                        \tag{C.EL.4}
\]
These contractions also exist on a compact total threefold over a
compact curve. Iteration on eligible leaves terminates, retaining
the punctured torus family, with exact Euler comparison
\[
 e(S)=e(S_{\rm rem})+2a+\sum_{j=1}^{b}\bigl(2-2g(B_j)\bigr),
                                                        \tag{C.EL.5}
\]
where $a$ is the number of recorded point contractions and the
$B_j$ are the recorded curve centres. Reversing those actual
ordinary blowups restores the original marked support and all
its neighbourhood maps.

For every specified elliptic curve $B$, both the split and the
nonzero-extension pairs in (C.EL.3) occur in a compact torus
degeneration over $\PP^1$ with exactly one singular fibre. That
fibre has two components: a two-torus and the prescribed ruled
pair. Its Euler number is zero, and every other fibre is a smooth
two-torus.
\end{statement}

\begin{proof}
\par\medskip\noindent\textbf{The actual canonical and normal restrictions.}
A local generator of the rank-one direct image of the canonical
bundle evaluates to a holomorphic three-form on $U$, using the
disc coordinate to trivialize $K_\Delta$. Its divisor is vertical:
on every smooth torus fibre the corresponding nonzero canonical
form has no zeros. Thus, after shrinking the disc, its divisor
is $\sum_v k_vT_v$. Adding a common integer to the $k_v$ does not
change the following comparisons. Since $S$ is the full reduced
principal fibre, restriction to $T$ gives actual line isomorphisms
\[
 \OO_U(T)|_T\simeq\OO_T(-C),\qquad
 K_T\simeq\OO_T\bigl((k_P-k_T-1)C\bigr).
                                                        \tag{C.EL.6}
\]
The second map is adjunction applied to the evaluation section
after dividing by a local equation of $T$ to its order $k_T$.
It also specifies a meromorphic canonical section with exactly
the indicated divisor. At $C$ the reduced equation is $xy=t$;
its product trivialization gives
$N_{C/T}\otimes N_{C/P}\simeq\OO_C$.

For a ruling fibre $L\simeq\PP^1$, put $d=C\cdot L$ and
$a=k_P-k_T-1$. Intersecting (C.EL.6) with $L$ gives
$-2=ad$. Thus $d>0$, $a<0$, and the integer pairs are exactly
$(d,a)=(1,-2)$ and $(2,-1)$. The map $C\to B$ is finite:
it is a nonconstant proper map of smooth compact curves. In the
degree-one case it is an isomorphism and hence a section.
Adjunction along this section gives
$K_C=\OO_C(-C)$, proving both normal identifications in
(C.EL.1), including their actual dual-normal map.

In the degree-two case $K_T+C$ is trivial, so adjunction gives
$K_C\simeq\OO_C$ and $g(C)=1$. Riemann--Hurwitz restricts
$g(B)$ to zero or one. We exclude the latter by the actual
cohomology maps. The projection satisfies
$q_*\OO_T=\OO_B$ and $R^1q_*\OO_T=0$, by the cohomology of
$\PP^1$ in a local bundle chart. The restriction
\[
 H^1(T,\OO_T)\longrightarrow H^1(C,\OO_C)
                                                        \tag{C.EL.7}
\]
is, under this Leray isomorphism, pullback by $C\to B$.
For $g(B)=1$ that degree-two map is etale; its trace composed
with pullback is multiplication by two. Hence (C.EL.7) is
nonzero. On the other hand the residue exact sequence
\[
 0\longrightarrow K_T\longrightarrow K_T(C)
       \longrightarrow K_C\longrightarrow0
                                                        \tag{C.EL.8}
\]
has $H^0(K_T)=0$, $H^0(K_T(C))=\C$ and $H^0(K_C)=\C$.
Its residue map is consequently an isomorphism and its
connecting map is zero. Serre duality identifies that connecting
map with the dual of (C.EL.7), by the residue pairing. This
contradiction excludes $g(B)=1$. Thus $B=\PP^1$ and
$T=\mathbb F_e$. The negative section $C_0$ has
$(-K_T)\cdot C_0=2-e$. A distinct irreducible effective
anticanonical curve has nonnegative intersection with $C_0$,
so $e\le2$. This proves (C.EL.2).

On the plane, write $C\sim dH$. Equation (C.EL.6) instead
gives $ad=-3$. Thus $(d,a)$ is $(1,-3)$ or $(3,-1)$,
giving exactly the stated line and cubic alternatives.

\par\medskip\noindent\textbf{The corrected analytic contraction theorem and its application.}
We use the Main Theorem of \cite[p.~484]{Nakano1970}, in its
corrected and strengthened form in
\cite[\S2, Theorem~2' and Proposition~9']{FujikiNakano1971}.
For a smooth divisor which is a holomorphic projective-space
bundle, the condition that its ambient normal line restricts to
$\OO(-1)$ on every projective fibre produces a smooth analytic
blowdown with that bundle as its exceptional divisor. The
supplement repairs the cohomology argument and proves that this
normal condition suffices. These are the versions read and used
here; their local blowdown maps glue by the actual isomorphisms
off the exceptional divisor.

In the section case, (C.EL.6) restricts to $\OO_L(-1)$ on
every ruling fibre. The theorem therefore supplies the stated
proper holomorphic blowdown $\pi:U\to U'$ along the given
bundle projection. It is the ordinary blowup of a smooth curve
$B$ in a smooth threefold. The equality
$\pi_*\OO_U=\OO_{U'}$ gives a unique holomorphic $f'$ with
$f=f'\circ\pi$. This equality can also be checked in the two
ordinary blowup charts by extending holomorphic functions
across the smooth codimension-two centre. Properness follows
from $f'^{-1}(K)=\pi(f^{-1}(K))$ for every compact
$K\subset\Delta$.

The coefficient of $T$ in $f^{-1}(0)$ is one, so the first
normal term of $f'$ along $B$ is a section of $N_{B/U'}^*$.
Its zero set on $\PP_B(N_{B/U'})$ is exactly $C$.
It is nonzero at every point of $B$: a zero at $b$ would
make that zero set contain the entire ruling fibre over $b$,
whereas $C$ is a section. Its normal differential is therefore
surjective along $B$. Thus $f'$ is a submersion there and
$P'=f'^{-1}(0)$ is smooth near the centre. The restriction of
the blowup to $P'$ is the blowup of its Cartier curve $B$,
which is an isomorphism. This proves $P\simeq P'$ with
the specified curve identification. Every other component
misses the exceptional divisor, so no other central component
meets this centre. The new full fibre remains reduced and SNC.

The normal exact sequence of $B\subset P'\subset U'$ is
precisely the first row of (C.EL.3): its quotient is trivialized
by $df'$, and its subline is $N_{C/P}\simeq K_B$ from
(C.EL.1). The exceptional divisor of the ordinary blowup is
the lines projectivization of this very normal bundle. Its
intersection with the strict central component is
$\PP_B(K_B)$. The quotient $V\to\OO_B$ defines this section
as the zero divisor of the corresponding section of
$\OO_{\PP_B(V)}(1)$, which gives the second exact line
identification in (C.EL.3). Serre duality gives
$H^1(B,K_B)\simeq H^0(B,\OO_B)^*\simeq\C$.
In local normal frames the transition matrices have the form
\[
        \begin{pmatrix}k_{ij}&a_{ij}\\0&1\end{pmatrix}.
                                                        \tag{C.EL.9}
\]
Changing a lift of the quotient changes $a_{ij}$ by its
actual Cech coboundary; rescaling the subline multiplies the
extension class by that scalar. Thus there are the zero and
nonzero orbits when endpoint frames are allowed to change.
The contraction retains the matrices (C.EL.9), rather than
discarding this extension class.

In the plane-line case the normal line is $\OO_{\PP^2}(-1)$.
The same corrected theorem, with centre a point, produces an
ordinary point blowdown onto a smooth threefold. The exceptional
coefficient one says that the descended $f'$ has a nonzero
linear term at that point. It is consequently a submersion
there. Its central surface pulls back to its point blowup,
with exceptional curve the specified line $C$, proving
(C.EL.4). The normal-dual map gives $C^2_P=-1$, consistently
with this actual surface blowdown.

\par\medskip\noindent\textbf{Euler and neighbourhood comparisons.}
For a curve contraction, the map is an isomorphism off $T$,
with fibre $\PP^1$ over each point of $B$. Additivity gives
$e(S)-e(S')=e(T)-e(B)=e(B)=2-2g(B)$.
For a point contraction it gives $e(T)-1=3-1=2$.
The same maps on a compact total space give the identical
Euler changes there. The descended family is unchanged over
the punctured base, and its fibres there are the original
marked two-tori. Each contraction removes one central
component; every iteration therefore terminates. Summing
these equalities gives (C.EL.5). The record consists of the
actual smooth centre, its normal extension, the local ordinary
blowup charts and their unchanged exterior identifications.
Reversing the maps restores the initial threefold near its
support, so this comparison retains all neighbourhood data.

\par\medskip\noindent\textbf{Both elliptic normal extensions occur compactly.}
Fix the specified elliptic $B$, choose another elliptic curve
$E$, and set $A=B\times E$. Use the compact smooth twistor
torus family (C.PT.13) with period order $R=1$ or $R=2$.
Its central torus is exactly $A$. Fix the curve $B\times\{0\}$
in that central member. In the actual primitive period cover
near this curve, (C.PR.13) writes the deck transformation as
\[
 (p,w,t)\longmapsto
       (p+\delta(t),\gamma(t)w,t),\qquad
 \delta(t)=t^R d(t),\quad d(0)\ne0,
                                                        \tag{C.EL.10}
\]
where $B=\C^*/\gamma(0)^\Z$ and $0<|\gamma(0)|<1$.
The curve is $p=t=0$. Differentiating these actual deck
coordinates gives its normal transition over this elliptic
cover:
\[
 V:\quad
 \begin{pmatrix}p\\t\end{pmatrix}\longmapsto
 \begin{pmatrix}1&\delta'(0)\\0&1\end{pmatrix}
 \begin{pmatrix}p\\t\end{pmatrix}.
                                                        \tag{C.EL.11}
\]
The term $\gamma'(0)w$ is tangent to the central elliptic
curve and hence vanishes in this normal quotient. The other
primitive period has identity normal transition.

For $R=2$, (C.EL.11) is the identity, so the normal extension
is split. For $R=1$, its upper entry is nonzero. A splitting
would supply a holomorphic $b$ on $\C^*$ satisfying
$b(\gamma(0)w)-b(w)=\delta'(0)$, with one possible opposite
sign from changing the frame convention. The Laurent
expansion on $\C^*$ has zero constant coefficient on the
left and a nonzero one on the right. No such splitting
exists. Thus its extension class is nonzero in $H^1(B,\OO_B)$.
Rescaling the transverse frame realizes every nonzero class
with its prescribed endpoint identifications.

Blow up this actual central curve in the compact smooth
twistor family. The exceptional marked pair is exactly
$\bigl(\PP_B(V),\PP_B(\OO_B)\bigr)$, by the ordinary
normal-bundle description just proved. The strict central
torus is unchanged because its centre is Cartier there.
Ordinary coordinates for the ideal $(p,t)$ have $t=xy$,
so the full new fibre has its two components with
multiplicity one. Its Euler number is zero by (C.EL.3),
and all other fibres are the unchanged smooth two-tori.

The exact degree alternatives, line-frame comparisons,
ordinary blowup Jacobians, normal-extension matrices and
Euler identities are computed in
\path{code/root_exceptional_leaves.py}.
\end{proof}
